% GENERATED FILE. Edit canonical lessons, then run npm run generate:books.
% canonical-source-hash: fnv1a64:12b1bbd184b75480
% canonical-lessons: FR-C477-cuisse, FR-C477-mollet, FR-C477-talon, FR-C477-orteil, FR-C477-hanche

\chapter{De la hanche au talon}
\label{ch:fr-de-la-hanche-au-talon}
\hlchaptermodality{477}

\begin{chapteropening}
\textbf{By the end of this chapter:} \emph{I can name the parts of the leg and foot.}

Complete the last lesson of chapter 477: I can name the parts of the leg and foot.
\end{chapteropening}

\section[la cuisse]{la cuisse — thigh}

\label{lesson:FR-C477-cuisse}



\begin{quote}
Before the new one: say the French for optimistic, then the French for pessimistic.
\end{quote}

\subsection*{You'll want to know: la cuisse}
\textbf{la cuisse} — \textquotedblleft{}thigh\textquotedblright{}.

\begin{quote}
Après la course, j'ai mal aux cuisses.
\end{quote}

\begin{quote}
\emph{After the run my thighs hurt.}
\end{quote}

Also a leg of chicken: \emph{une cuisse de poulet}.

\begin{cousinweb}
\emph{la cuisse} is inherited, through Old French, from Latin \emph{coxa} \textquotedblleft{}hip\textquotedblright{} (Lewis \& Short); the word slid down from the hip to the thigh.
\end{cousinweb}

\begin{grammarlens}[title={What you've built}]
One more word for telling, explaining and describing at B1.
\end{grammarlens}

\subsection*{Your turn}
\begin{itemize}
\raggedright
  \item \emph{Say it:} \emph{la cuisse}
  \item \emph{Say it:} \emph{la cuisse}, once more
  \item \emph{Say it:} a sentence of your own with \emph{la cuisse}
  \item \emph{Recall:} say the French for optimistic, then the French for pessimistic, then say \emph{la cuisse} again
\end{itemize}

\begin{tcolorbox}[breakable,colback=teal!4,colframe=teal!35!black,title={Before you move on}]
What does la cuisse mean? (\textquotedblleft{}Thigh\textquotedblright{}.) Say it once more.
\end{tcolorbox}
\section[le mollet]{le mollet — calf (of the leg)}

\label{lesson:FR-C477-mollet}



\begin{quote}
Before the new one: say the French for pessimistic, then the French for thigh.
\end{quote}

\subsection*{You'll want to know: le mollet}
\textbf{le mollet} — \textquotedblleft{}calf (of the leg)\textquotedblright{}.

\begin{quote}
Le cycliste a des mollets très forts.
\end{quote}

\begin{quote}
\emph{The cyclist has very strong calves.}
\end{quote}

The soft back part of the lower leg.

\begin{cousinweb}
\emph{le mollet} is built on \emph{mol}, the old form of \emph{mou} \textquotedblleft{}soft\textquotedblright{}, with the small ending \emph{-et} (EtymDB); \emph{mou} is inherited from Latin \emph{mollis} \textquotedblleft{}soft\textquotedblright{} (Lewis \& Short).
\end{cousinweb}

\begin{grammarlens}[title={What you've built}]
One more word for telling, explaining and describing at B1.
\end{grammarlens}

\subsection*{Your turn}
\begin{itemize}
\raggedright
  \item \emph{Say it:} \emph{le mollet}
  \item \emph{Say it:} \emph{le mollet}, once more
  \item \emph{Say it:} a sentence of your own with \emph{le mollet}
  \item \emph{Recall:} say the French for pessimistic, then the French for thigh, then say \emph{le mollet} again
\end{itemize}

\begin{tcolorbox}[breakable,colback=teal!4,colframe=teal!35!black,title={Before you move on}]
What does le mollet mean? (\textquotedblleft{}Calf (of the leg)\textquotedblright{}.) Say it once more.
\end{tcolorbox}
\section[le talon]{le talon — heel}

\label{lesson:FR-C477-talon}



\begin{quote}
Before the new one: say the French for thigh, then the French for calf.
\end{quote}

\subsection*{You'll want to know: le talon}
\textbf{le talon} — \textquotedblleft{}heel\textquotedblright{}.

\begin{quote}
Ces chaussures me font mal au talon.
\end{quote}

\begin{quote}
\emph{These shoes hurt my heel.}
\end{quote}

Also the heel of a shoe: \emph{des talons hauts} — high heels.

\begin{cousinweb}
\emph{le talon} is inherited from Vulgar Latin \emph{tālo} (EtymDB), from Latin \emph{tālus} \textquotedblleft{}the ankle, the anklebone\textquotedblright{} (Lewis \& Short).
\end{cousinweb}

\begin{grammarlens}[title={What you've built}]
One more word for telling, explaining and describing at B1.
\end{grammarlens}

\subsection*{Your turn}
\begin{itemize}
\raggedright
  \item \emph{Say it:} \emph{le talon}
  \item \emph{Say it:} \emph{le talon}, once more
  \item \emph{Say it:} a sentence of your own with \emph{le talon}
  \item \emph{Recall:} say the French for thigh, then the French for calf, then say \emph{le talon} again
\end{itemize}

\begin{tcolorbox}[breakable,colback=teal!4,colframe=teal!35!black,title={Before you move on}]
What does le talon mean? (\textquotedblleft{}Heel\textquotedblright{}.) Say it once more.
\end{tcolorbox}
\section[l'orteil]{l'orteil — toe}

\label{lesson:FR-C477-orteil}



\begin{quote}
Before the new one: say the French for calf, then the French for heel.
\end{quote}

\subsection*{You'll want to know: l'orteil}
\textbf{l'orteil} — \textquotedblleft{}toe\textquotedblright{}.

\begin{quote}
Je me suis cogné l'orteil contre un pied de table.
\end{quote}

\begin{quote}
\emph{I stubbed my toe on the leg of a table.}
\end{quote}

Masculine. \emph{Le gros orteil} — the big toe. Fingers are \emph{les doigts}; toes can also be \emph{les doigts de pied}.

\begin{cousinweb}
\emph{l'orteil} comes from Old French \emph{arteil}, from Latin \emph{articulus} \textquotedblleft{}a small joint, a limb\textquotedblright{} (EtymDB; Lewis \& Short); its first vowel was later changed.
\end{cousinweb}

\begin{grammarlens}[title={What you've built}]
One more word for telling, explaining and describing at B1.
\end{grammarlens}

\subsection*{Your turn}
\begin{itemize}
\raggedright
  \item \emph{Say it:} \emph{l'orteil}
  \item \emph{Say it:} \emph{l'orteil}, once more
  \item \emph{Say it:} a sentence of your own with \emph{l'orteil}
  \item \emph{Recall:} say the French for calf, then the French for heel, then say \emph{l'orteil} again
\end{itemize}

\begin{tcolorbox}[breakable,colback=teal!4,colframe=teal!35!black,title={Before you move on}]
What does l'orteil mean? (\textquotedblleft{}Toe\textquotedblright{}.) Say it once more.
\end{tcolorbox}
\section[la hanche]{la hanche — hip}

\label{lesson:FR-C477-hanche}



\begin{quote}
Before the new one: say the French for heel, then the French for toe.
\end{quote}

\subsection*{You'll want to know: la hanche}
\textbf{la hanche} — \textquotedblleft{}hip\textquotedblright{}.

\begin{quote}
Elle porte son bébé sur la hanche.
\end{quote}

\begin{quote}
\emph{She carries her baby on her hip.}
\end{quote}

The \emph{h} is aspirated: \emph{la hanche}, not \emph{l'hanche}.

\begin{cousinweb}
\emph{la hanche} comes from a Germanic word for the haunch. English \emph{haunch} is the same word, taken from Old French (EtymDB).
\end{cousinweb}

\begin{grammarlens}[title={What you've built}]
One more word for telling, explaining and describing at B1.
\end{grammarlens}

\subsection*{Your turn}
\begin{itemize}
\raggedright
  \item \emph{Say it:} \emph{la hanche}
  \item \emph{Say it:} \emph{la hanche}, once more
  \item \emph{Say it:} a sentence of your own with \emph{la hanche}
  \item \emph{Recall:} say the French for heel, then the French for toe, then say \emph{la hanche} again
\end{itemize}

\begin{tcolorbox}[breakable,colback=teal!4,colframe=teal!35!black,title={Before you move on}]
What does la hanche mean? (\textquotedblleft{}Hip\textquotedblright{}.) Say it once more.
\end{tcolorbox}
